\chapter{非线性控制}

\intro{
前面各章的研究对象是线性系统，其理论基础建立在叠加原理之上。然而，绝大多数实际物理系统本质上是非线性的——从机械臂的多轴耦合到飞机的高攻角气动特性，从化学反应的指数动力学到电力系统的频率振荡。非线性系统呈现出线性系统不具备的丰富动力学行为：多平衡点、极限环、分岔和混沌。本章系统介绍非线性控制的基本理论和方法，包括Lyapunov稳定性理论、滑模控制、反馈线性化、Backstepping和非线性观测器设计。
}

%%%%%%%%%%%%%%%%%%%%%%%%%%%%%%%%%%%%%%%%%%%%%%%%%%%%%%%%%%%%%%%%%%%%%%%%%%%%%%%%
\section{非线性系统的特性}

\subsection{非线性系统的基本描述}

一般形式的非线性系统可写为
\[
\dot{\mathbf{x}} = \mathbf{f}(\mathbf{x}, \mathbf{u}, t),\qquad \mathbf{y} = \mathbf{h}(\mathbf{x}, \mathbf{u}, t)
\]
其中$\mathbf{f}$和$\mathbf{h}$是非线性向量函数。当$\mathbf{f}$不显含时间$t$时，系统称为\textbf{自治系统}：
\[
\dot{\mathbf{x}} = \mathbf{f}(\mathbf{x})
\]

仿射非线性系统是控制理论中最常研究的一类：
\begin{myformula}
\dot{\mathbf{x}} = \mathbf{f}(\mathbf{x}) + \sum_{i=1}^{m} \mathbf{g}_i(\mathbf{x})u_i = \mathbf{f}(\mathbf{x}) + G(\mathbf{x})\mathbf{u}
\end{myformula}
其中$\mathbf{f}(\mathbf{x})$是漂移向量场，$\mathbf{g}_i(\mathbf{x})$是控制向量场。

\subsection{叠加原理的失效与多平衡点}

线性系统$\dot{\mathbf{x}} = A\mathbf{x}$的平衡点只有原点（当$A$非奇异时）。而非线性系统$\dot{\mathbf{x}} = \mathbf{f}(\mathbf{x})$可能有多个孤立平衡点，每个平衡点周围的动力学行为可能完全不同。例如，单摆方程
\[
\ddot{\theta} + \frac{g}{\ell}\sin\theta = 0
\]
具有无穷多个平衡点$\theta = k\pi\;(k \in \mathbb{Z})$：偶数倍$\pi$处是稳定焦点，奇数倍$\pi$处是鞍点。

\begin{remark}
叠加原理的失效意味着：不能将非线性系统的响应分解为不同输入分别响应的叠加，也不能将复杂运动分解为模态的组合。这是非线性控制比线性控制更具挑战性的根本原因。
\end{remark}

\subsection{极限环、混沌与分岔}

\textbf{极限环}是非线性系统独有的孤立周期轨道，在线性系统中不可能出现。典型的例子是Van der Pol振荡器：
\[
\ddot{x} + \mu(x^2 - 1)\dot{x} + x = 0,\quad \mu > 0
\]
对于任何非零初始条件，解轨线都收敛到同一个极限环。

\textbf{混沌}表现为对初值的极端敏感性（"蝴蝶效应"）、非周期长期行为和奇怪吸引子（Strange Attractor）。Lorenz系统是最经典的混沌示例：
\begin{align*}
\dot{x} &= \sigma(y - x) \\
\dot{y} &= x(\rho - z) - y \\
\dot{z} &= xy - \beta z
\end{align*}

\textbf{分岔}（Bifurcation）是指系统参数穿过临界值时，相图拓扑结构发生的定性改变。常见类型包括鞍结分岔（saddle-node）、Hopf分岔和叉式分岔（pitchfork）。

\begin{figure}[H]
\centering
\begin{tikzpicture}[nlCharStyle/.style={thick}, nlCharScale/.style={scale=1.1}]
    % Stable node
    \begin{scope}[nlCharScale]
        \draw[nlCharStyle, ->] (-2,0) -- (2,0) node[right] {$x_1$};
        \draw[nlCharStyle, ->] (0,-2) -- (0,2) node[above] {$x_2$};
        % Trajectories converging to origin
        \foreach \ang in {0, 30, 60, 90, 120, 150, 180, 210, 240, 270, 300, 330} {
            \draw[blue!60, ->] ({1.8*cos(\ang)}, {1.8*sin(\ang)}) -- ({0.3*cos(\ang)}, {0.3*sin(\ang)});
        }
        \node at (2.5, 1.5) {\small 稳定节点};
    \end{scope}
    
    % Stable focus
    \begin{scope}[nlCharScale, xshift=6cm]
        \draw[nlCharStyle, ->] (-2,0) -- (2,0) node[right] {$x_1$};
        \draw[nlCharStyle, ->] (0,-2) -- (0,2) node[above] {$x_2$};
        \draw[red!70, ->, domain=0:6*pi, samples=200, variable=\t]
            plot ({exp(-0.25*\t)*cos(\t r)}, {exp(-0.25*\t)*sin(\t r)});
        \node at (2.5, 1.5) {\small 稳定焦点};
    \end{scope}
    
    % Saddle
    \begin{scope}[nlCharScale, xshift=12cm]
        \draw[nlCharStyle, ->] (-2,0) -- (2,0) node[right] {$x_1$};
        \draw[nlCharStyle, ->] (0,-2) -- (0,2) node[above] {$x_2$};
        \draw[green!60!black, ->] (-1.8, 0.05) -- (-0.15, 0.02);
        \draw[green!60!black, ->] (1.8, -0.05) -- (0.15, -0.02);
        \draw[green!60!black, ->] (-0.05, 1.8) -- (-0.02, 0.15);
        \draw[green!60!black, ->] (0.05, -1.8) -- (0.02, -0.15);
        \draw[green!60!black, dashed] (-1.5, -1.5) -- (0, 0);
        \draw[green!60!black, dashed] (1.5, 1.5) -- (0, 0);
        \draw[green!60!black, dashed] (-1.5, 1.5) -- (0, 0);
        \draw[green!60!black, dashed] (1.5, -1.5) -- (0, 0);
        \node at (2.5, 1.5) {\small 鞍点};
    \end{scope}
\end{tikzpicture}
\caption{典型相平面轨迹类型。左：稳定节点（特征值为负实根），中：稳定焦点（特征值为共轭复根），右：鞍点（特征值为一正一负实根）。}
\label{fig:ch11_phase_portrait}
\end{figure}

\begin{figure}[H]
\centering
\begin{tikzpicture}[limCycleStyle/.style={thick}]
    \draw[limCycleStyle, ->] (-2.5,0) -- (2.5,0) node[right] {$x_1$};
    \draw[limCycleStyle, ->] (0,-2.5) -- (0,2.5) node[above] {$x_2$};
    % Limit cycle ellipse
    \draw[red, very thick] (0,0) ellipse (1.8 and 1.2);
    % Inward trajectory
    \draw[blue!70, ->, domain=0:4*pi, samples=200, variable=\t]
        plot ({2.5*exp(-0.1*\t)*cos(\t r)}, {1.8*exp(-0.1*\t)*sin(\t r)});
    % Outward trajectory
    \draw[green!50!black, ->, domain=0:3*pi, samples=200, variable=\t]
        plot ({0.3*exp(0.15*\t)*cos(\t r)}, {0.2*exp(0.15*\t)*sin(\t r)});
    \node at (0, -2.8) {\small 稳定极限环（内外轨线均收敛至环）};
\end{tikzpicture}
\caption{极限环示意图。外部轨线向内螺旋收敛，内部轨线向外螺旋膨胀，最终都趋向红色椭圆所示的极限环。}
\label{fig:ch11_limit_cycle}
\end{figure}

%%%%%%%%%%%%%%%%%%%%%%%%%%%%%%%%%%%%%%%%%%%%%%%%%%%%%%%%%%%%%%%%%%%%%%%%%%%%%%%%
\section{Lyapunov稳定性理论}

\subsection{自治系统的稳定性概念}

考虑自治系统$\dot{\mathbf{x}} = \mathbf{f}(\mathbf{x})$，设$\mathbf{x} = \mathbf{0}$是一个平衡点（通过坐标平移总可达到），即$\mathbf{f}(\mathbf{0}) = \mathbf{0}$。

\begin{mydefinition}{Lyapunov意义下的稳定性}
\begin{enumerate}
\item \textbf{稳定}：$\forall \varepsilon > 0$，$\exists \delta(\varepsilon) > 0$使得$\|\mathbf{x}(0)\| < \delta \Rightarrow \|\mathbf{x}(t)\| < \varepsilon$，$\forall t \ge 0$；
\item \textbf{渐近稳定}：稳定，且$\lim_{t\to\infty}\|\mathbf{x}(t)\| = 0$；
\item \textbf{指数稳定}：$\exists \alpha, \beta > 0$使得$\|\mathbf{x}(t)\| \le \beta\|\mathbf{x}(0)\|e^{-\alpha t}$；
\item \textbf{全局渐近稳定}：渐近稳定，且对任意初始条件$\mathbf{x}(0) \in \R^n$均收敛到原点。
\end{enumerate}
\end{mydefinition}

\begin{remark}
与线性系统不同，非线性系统的稳定性往往是\textbf{局部}性质——一个平衡点可能在某些初始条件下稳定，而在另一些初始条件下不稳定。因此非线性稳定性理论中，"吸引域"（Region of Attraction）的概念至关重要。
\end{remark}

\subsection{Lyapunov直接法（能量法）}

Lyapunov直接法的核心思想是：如果能够找到一个"能量函数"$V(\mathbf{x})$，它沿系统轨线的导数始终非正（表示能量不断衰减），那么系统的平衡点就是稳定的。

\begin{myTheorem}{Lyapunov稳定性定理}
设$\mathbf{x} = \mathbf{0}$是系统$\dot{\mathbf{x}} = \mathbf{f}(\mathbf{x})$的一个平衡点。若存在连续可微函数$V: D \to \R$（$D \subset \R^n$包含原点）满足
\begin{enumerate}
\item $V(\mathbf{0}) = 0$，且$V(\mathbf{x}) > 0$对任意$\mathbf{x} \in D \setminus \{\mathbf{0}\}$（正定性）；
\item $\dot{V}(\mathbf{x}) = \Nabla V(\mathbf{x}) \cdot \mathbf{f}(\mathbf{x}) \le 0$对任意$\mathbf{x} \in D$（半负定性）。
\end{enumerate}
则原点是\textbf{稳定}的。若进一步$\dot{V}(\mathbf{x}) < 0$对任意$\mathbf{x} \in D \setminus \{\mathbf{0}\}$，则原点是\textbf{渐近稳定}的。
\end{myTheorem}

\begin{proof}[Lyapunov稳定性定理的简要证明思路]
正定性保证$V$的等值面构成一族围绕原点的闭合曲面。$\dot{V} \le 0$保证轨线不会穿越等值面向外发展——轨线始终停留在初始$V$值对应的等值面包围的区域内。通过选取充分小的初始$V$值，可使轨线始终保持在任意给定的$\varepsilon$邻域内。
\end{proof}

\subsection{Lyapunov函数的构造方法}

构造合适的Lyapunov函数是非线性稳定性分析中的主要挑战。常用方法包括：

\begin{enumerate}
\item \textbf{能量法}：对于物理系统，总能量（动能+势能）往往是天然的Lyapunov函数候选。例如，对于质量-弹簧-阻尼系统，总机械能$V = \frac{1}{2}m\dot{x}^2 + \frac{1}{2}kx^2$的导数为$\dot{V} = -c\dot{x}^2 \le 0$；
\item \textbf{平方和法}：对于二次型$V = \mathbf{x}^{\top}P\mathbf{x}$，$\dot{V} = \mathbf{x}^{\top}(A^{\top}P + PA)\mathbf{x}$，这对应于线性情况下的Lyapunov方程；
\item \textbf{Krasovskii方法}：取$V = \mathbf{f}^{\top}P\mathbf{f}$，适用于某些特殊结构；
\item \textbf{反推法（Backstepping）}：对级联系统逐级构造Lyapunov函数（参见第5节）。
\end{enumerate}

\begin{example}[单摆的Lyapunov分析]
考虑带阻尼的单摆方程：$\ddot{\theta} + b\dot{\theta} + \frac{g}{\ell}\sin\theta = 0$（$b > 0$）。取状态变量$x_1 = \theta$，$x_2 = \dot{\theta}$，则状态方程：
\[
\dot{x}_1 = x_2,\quad \dot{x}_2 = -\frac{g}{\ell}\sin x_1 - b x_2
\]
取Lyapunov函数为总能量：
\[
V = \frac{1}{2}x_2^2 + \frac{g}{\ell}(1 - \cos x_1)
\]
求导：
\[
\dot{V} = x_2\dot{x}_2 + \frac{g}{\ell}\sin x_1 \cdot \dot{x}_1
= x_2\left(-\frac{g}{\ell}\sin x_1 - b x_2\right) + \frac{g}{\ell}\sin x_1 \cdot x_2 = -b x_2^2 \le 0
\]
由Lyapunov定理，原点$(0, 0)$是稳定的。再利用Lasalle不变原理可以进一步证明渐近稳定性。
\end{example}

\subsection{Lasalle不变原理}

Lasalle不变原理推广了Lyapunov定理，使我们即使在$\dot{V} \le 0$（非严格负定）时也能推断渐近稳定性。

\begin{myTheorem}{Lasalle不变原理}
设$\Omega \subset D$是一个紧致正不变集（即系统轨线一旦进入$\Omega$就永不离开）。设$V: D \to \R$连续可微，在$\Omega$内满足$\dot{V} \le 0$。令$E = \{\mathbf{x} \in \Omega \mid \dot{V}(\mathbf{x}) = 0\}$，$M$是$E$中的最大不变集。则当$t \to \infty$时，任何起始于$\Omega$的轨线都趋于$M$。
\end{myTheorem}

\begin{figure}[H]
\centering
\begin{tikzpicture}[lyapStyle/.style={thick}]
    % Coordinate axes
    \draw[lyapStyle, ->] (-2.5, 0) -- (2.5, 0) node[right] {$x_1$};
    \draw[lyapStyle, ->] (0, -2.5) -- (0, 2.5) node[above] {$x_2$};
    
    % Level sets of V
    \foreach \r in {0.5, 1.0, 1.5, 2.0} {
        \draw[blue!40, dashed] (0, 0) ellipse ({\r*1.0} and {\r*0.7});
    }
    
    % Trajectory going inward
    \draw[red, ->, domain=0:8*pi, samples=200, variable=\t]
        plot ({2.2*exp(-0.15*\t)*cos(\t r)}, {1.5*exp(-0.15*\t)*sin(\t r)});
    
    \node[blue!70] at (1.7, 1.2) {$V = c_1$};
    \node[blue!70] at (2.1, 1.5) {$V = c_2 > c_1$};
    
    \node at (0, -3) {\small 轨线不断穿越$V$的等值面（向内的方向），$\dot{V} \le 0$};
\end{tikzpicture}
\caption{Lyapunov函数的等值面与系统轨线的关系。轨线从外向内穿越递减的$V$等值面，体现了"能量"不断耗散的过程。}
\label{fig:ch11_lyapunov_levelsets}
\end{figure}

%%%%%%%%%%%%%%%%%%%%%%%%%%%%%%%%%%%%%%%%%%%%%%%%%%%%%%%%%%%%%%%%%%%%%%%%%%%%%%%%
\section{滑模控制}

\subsection{滑模控制的基本原理}

滑模控制（Sliding Mode Control, SMC）是一种鲁棒非线性控制方法，其核心思路是：在状态空间中设计一个"滑模面"$s(\mathbf{x}) = 0$，然后用不连续的控制律将系统状态"驱动"到该滑模面上，并使其沿滑模面滑向平衡点。

对于单输入仿射非线性系统：
\[
\dot{\mathbf{x}} = \mathbf{f}(\mathbf{x}) + \mathbf{g}(\mathbf{x})u
\]
设计滑模面$s(\mathbf{x})$，使$s = 0$是一个$(n-1)$维超曲面，且在该面上的"滑模动态"是渐近稳定的。

\subsection{滑模面的设计}

常用的滑模面设计方法包括：

\textbf{(1) 线性滑模面}：$s = \mathbf{c}^{\top}\mathbf{x}$，其中$\mathbf{c} \in \R^n$是待设计的常数向量。这种方式等价于将滑模运动配置为$n-1$维线性系统的稳定运动。对于可化为正则形式的系统，$\mathbf{c}$的选择对应着滑模极点配置。

\textbf{(2) 终端滑模面}：引入分数幂项以实现在有限时间内收敛：
\[
s = \dot{x} + \beta x^{q/p},\quad p > q > 0,\;p, q\text{为奇数}
\]

\textbf{(3) 积分滑模面}：引入误差的积分项以消除稳态误差：
\[
s = \dot{e} + \lambda_1 e + \lambda_2 \int_0^t e(\tau)d\tau
\]

\subsection{趋近律设计}

滑模控制的关键是设计控制律使$s \to 0$（到达阶段）并保持$s = 0$（滑动阶段）。趋近律$\dot{s}$的设计决定了收敛到滑模面的速度和质量。

\begin{mydefinition}{常用趋近律}
\begin{enumerate}
\item \textbf{等速趋近律}：$\dot{s} = -\eta\operatorname{sgn}(s)$，$\eta > 0$为趋近速度；
\item \textbf{指数趋近律}：$\dot{s} = -\eta\operatorname{sgn}(s) - ks$，$k, \eta > 0$。指数项$-ks$在远离滑模面时提供快速趋近，等速项$-\eta\operatorname{sgn}(s)$保证在$s=0$附近不会"滑出"；
\item \textbf{幂次趋近律}：$\dot{s} = -k|s|^{\alpha}\operatorname{sgn}(s)$，$k > 0$，$0 < \alpha < 1$；
\item \textbf{一般趋近律}：$\dot{s} = -\eta\operatorname{sgn}(s) - f(s)$，$f(0) = 0$且$sf(s) > 0$当$s \neq 0$。
\end{enumerate}
\end{mydefinition}

控制律的构造基于Lyapunov函数$V = \frac{1}{2}s^2$：
\[
\dot{V} = s\dot{s} = s(\Nabla s \cdot \dot{\mathbf{x}}) = s\left(\frac{\partial s}{\partial \mathbf{x}}\mathbf{f} + \frac{\partial s}{\partial \mathbf{x}}\mathbf{g}u\right)
\]
取控制律
\[
u = -\left(\frac{\partial s}{\partial \mathbf{x}}\mathbf{g}\right)^{-1}\left(\frac{\partial s}{\partial \mathbf{x}}\mathbf{f} + \eta\operatorname{sgn}(s)\right)
\]
则$\dot{V} = -\eta|s|$，保证有限时间内到达滑模面。

\subsection{抖振问题与边界层方法}

理想滑模控制要求控制器在$s=0$处瞬间切换符号（符号函数$\operatorname{sgn}(s)$），这在物理上不可能实现（执行器带宽有限），直接实施会产生高频振荡——称为\textbf{抖振}（Chattering）。

\textbf{边界层方法}将不连续的$\operatorname{sgn}$函数替换为饱和函数$\operatorname{sat}$以平滑控制信号：
\[
\operatorname{sat}(s/\Phi) = \begin{cases}
\operatorname{sgn}(s), & |s| > \Phi \\
s/\Phi, & |s| \le \Phi
\end{cases}
\]
其中$\Phi > 0$是边界层厚度。在边界层内使用连续控制消除了抖振，但付出的代价是滑模运动的理想鲁棒性在边界层内有所丧失（系统状态不再严格保持$s = 0$，而是保持在$|s| \le \Phi$内）。

\begin{figure}[H]
\centering
\begin{tikzpicture}[smcStyle/.style={thick}]
    \draw[smcStyle, ->] (-3, 0) -- (3, 0) node[right] {$x_1$};
    \draw[smcStyle, ->] (0, -2.5) -- (0, 2.5) node[above] {$x_2$};
    
    % Sliding surface
    \draw[red, ultra thick] (-2.5, -2.5) -- (2.5, 2.5);
    \node[red] at (2.8, 2.3) {$s = 0$};
    
    % Reaching phase
    \draw[blue!70, ->, domain=0:1.5, samples=100, variable=\t]
        plot ({-2.2 + 1.5*\t}, {1.8 - 1.2*\t});
    \node[blue!70] at (-1.5, 1.2) {到达阶段};
    
    % Sliding phase
    \draw[green!50!black, very thick, ->, domain=0:1.5, samples=50, variable=\t]
        plot ({-0.7 + 0.8*\t}, {-0.7 + 0.8*\t});
    \node[green!50!black] at (0.5, -0.3) {滑动阶段};
    
    % Chattering illustration near s=0
    \draw[purple, ->] (1.5, 1.5) -- (1.7, 1.4) -- (1.55, 1.65) -- (1.75, 1.55);
    \node[purple] at (3.2, 1.5) {抖振};
\end{tikzpicture}
\caption{滑模控制的阶段划分。蓝色轨迹为"到达阶段"，系统状态从初始位置被驱动到滑模面$s=0$上；绿色轨迹为"滑动阶段"，状态沿滑模面滑向原点。紫色放大部分示意抖振现象。}
\label{fig:ch11_sliding_mode}
\end{figure}

%%%%%%%%%%%%%%%%%%%%%%%%%%%%%%%%%%%%%%%%%%%%%%%%%%%%%%%%%%%%%%%%%%%%%%%%%%%%%%%%
\section{反馈线性化}

\subsection{输入-状态线性化}

反馈线性化（Feedback Linearization）的核心思想是通过坐标变换和非线性反馈将非线性系统\textbf{精确地}转化为线性系统（而非像Jacobian线性化那样只做近似）。考虑仿射非线性系统：
\[
\dot{\mathbf{x}} = \mathbf{f}(\mathbf{x}) + G(\mathbf{x})\mathbf{u}
\]

如果存在微分同胚$T: \R^n \to \R^n$（坐标变换$\mathbf{z} = T(\mathbf{x})$）和状态反馈
\[
\mathbf{u} = \boldsymbol{\alpha}(\mathbf{x}) + \boldsymbol{\beta}(\mathbf{x})\mathbf{v}
\]
使得在新坐标$\mathbf{z}$下系统变为线性可控形式（Brunovsky标准型）：
\[
\dot{\mathbf{z}} = A_c\mathbf{z} + B_c\mathbf{v}
\]
则称该系统\textbf{可输入-状态线性化}。

\subsection{Lie导数与Lie括号}

精确线性化的数学基础是微分几何中的Lie导数和Lie括号运算。

\begin{mydefinition}{Lie导数与Lie括号}
设$h: \R^n \to \R$是光滑标量函数，$\mathbf{f}: \R^n \to \R^n$是光滑向量场。$h$沿$\mathbf{f}$的\textbf{Lie导数}定义为
\[
L_{\mathbf{f}}h(\mathbf{x}) = \frac{\partial h}{\partial \mathbf{x}}\mathbf{f}(\mathbf{x}) = \sum_{i=1}^{n}\frac{\partial h}{\partial x_i}f_i(\mathbf{x})
\]
高阶Lie导数迭代定义：$L_{\mathbf{f}}^k h = L_{\mathbf{f}}(L_{\mathbf{f}}^{k-1}h)$。

两个向量场$\mathbf{f}$和$\mathbf{g}$的\textbf{Lie括号}（对易子）定义为
\[
[\mathbf{f}, \mathbf{g}](\mathbf{x}) = \frac{\partial \mathbf{g}}{\partial \mathbf{x}}\mathbf{f}(\mathbf{x}) - \frac{\partial \mathbf{f}}{\partial \mathbf{x}}\mathbf{g}(\mathbf{x})
\]
其中$\frac{\partial \mathbf{g}}{\partial \mathbf{x}}$和$\frac{\partial \mathbf{f}}{\partial \mathbf{x}}$分别是向量场$\mathbf{g}$和$\mathbf{f}$的Jacobian矩阵。
\end{mydefinition}

\subsection{相对阶与输入-输出线性化}

对于单输入单输出（SISO）非线性系统：
\[
\dot{\mathbf{x}} = \mathbf{f}(\mathbf{x}) + \mathbf{g}(\mathbf{x})u,\quad y = h(\mathbf{x})
\]
\textbf{相对阶}$r$定义为满足$L_{\mathbf{g}}L_{\mathbf{f}}^{r-1}h(\mathbf{x}) \neq 0$的最小整数。

\begin{example}[相对阶的计算]
对于系统$\ddot{y} + \sin y = u$，有
\begin{align*}
\dot{x}_1 &= x_2 = f_1 + g_1 u\\
\dot{x}_2 &= -\sin x_1 + u = f_2 + g_2 u\\
y &= x_1 = h(\mathbf{x})
\end{align*}
计算：$L_{\mathbf{g}}h = 0$，$L_{\mathbf{g}}L_{\mathbf{f}}h = L_{\mathbf{g}}(x_2) = 1 \neq 0$。因此相对阶$r = 2$。
\end{example}

若$r = n$（系统的阶数），则系统可完全的输入-状态线性化，没有零动态。若$r < n$，则存在$(n-r)$维的\textbf{零动态}（Zero Dynamics）——即当输出被强制为零时内部状态的动态行为。系统的内部稳定性取决于零动态的稳定性。

\begin{myTheorem}{零动态与最小相位性}
若系统的零动态是渐近稳定的，则称该系统为\textbf{最小相位}系统。输入-输出线性化控制器可以稳定整个系统当且仅当零动态是渐近稳定的。这与线性系统中"不稳定的零点限制可实现的闭环性能"具有深刻的类比关系。
\end{myTheorem}

\begin{remark}
反馈线性化与经典的Jacobian线性化有本质区别：前者是\textbf{精确的}代数变换，适用面更广但条件也更苛刻（需要验证对合性条件）；后者只是局部的一阶近似。在工程实践中，若精确线性化条件满足，应当优先使用反馈线性化以获得全局或大范围的有效性。
\end{remark}

%%%%%%%%%%%%%%%%%%%%%%%%%%%%%%%%%%%%%%%%%%%%%%%%%%%%%%%%%%%%%%%%%%%%%%%%%%%%%%%%
\section{Backstepping设计方法}

\subsection{Backstepping的基本思想}

Backstepping（反推法）是一种针对严格反馈（Strict-Feedback）系统的递归Lyapunov设计方法。其核心思路是将高阶系统分解为一系列级联的一阶子系统，从最内层的子系统开始，逐层向外"反推"设计虚拟控制律和Lyapunov函数。

\subsection{严格反馈形式}

严格反馈系统的标准形式为：
\begin{align*}
\dot{x}_1 &= f_1(x_1) + g_1(x_1)x_2 \\
\dot{x}_2 &= f_2(x_1, x_2) + g_2(x_1, x_2)x_3 \\
&\;\;\vdots \\
\dot{x}_{n-1} &= f_{n-1}(x_1, \dots, x_{n-1}) + g_{n-1}(x_1, \dots, x_{n-1})x_n \\
\dot{x}_n &= f_n(x_1, \dots, x_n) + g_n(x_1, \dots, x_n)u
\end{align*}
其中$g_i \neq 0$（能控性条件）。注意这种结构：每个状态方程只受下一状态的仿射影响，形成了"下三角"结构。

\subsection{二阶系统的Backstepping设计范例}

考虑二阶严格反馈系统：
\begin{align*}
\dot{x}_1 &= f_1(x_1) + g_1(x_1)x_2 \\
\dot{x}_2 &= f_2(x_1, x_2) + g_2(x_1, x_2)u
\end{align*}

\textbf{Step 1}：将$x_2$视为对子系统$x_1$的"虚拟控制"。选择期望的虚拟控制律$\alpha(x_1)$使得$x_1$子系统稳定。定义跟踪误差$z_1 = x_1$，$z_2 = x_2 - \alpha(x_1)$。

选取Lyapunov函数$V_1 = \frac{1}{2}z_1^2$，则
\[
\dot{V}_1 = z_1(f_1 + g_1(z_2 + \alpha)) = z_1(f_1 + g_1\alpha) + g_1 z_1 z_2
\]
取$\alpha = -\frac{1}{g_1}(f_1 + k_1 z_1)$（$k_1 > 0$），则$\dot{V}_1 = -k_1 z_1^2 + g_1 z_1 z_2$。

\textbf{Step 2}：构造增广的Lyapunov函数$V_2 = V_1 + \frac{1}{2}z_2^2$：
\begin{align*}
\dot{V}_2 &= -k_1 z_1^2 + g_1 z_1 z_2 + z_2(f_2 + g_2 u - \dot{\alpha}) \\
&= -k_1 z_1^2 + z_2(g_1 z_1 + f_2 + g_2 u - \dot{\alpha})
\end{align*}
取控制律
\[
u = \frac{1}{g_2}(-g_1 z_1 - f_2 + \dot{\alpha} - k_2 z_2)
\]
则$\dot{V}_2 = -k_1 z_1^2 - k_2 z_2^2 < 0$，系统全局渐近稳定。

\begin{remark}
Backstepping设计有两个显著优势：(1) 灵活性——在每个"反推"步骤中可以选择任意的Lyapunov函数$V_i$和控制Lyapunov函数（CLF），从而引入对非线性项的系统化补偿；(2) 避免了反馈线性化中"精确消去全部非线性"的要求，只需处理对稳定性不利的非线性项，有利于提高控制器的鲁棒性。
\end{remark}

\subsection{自适应Backstepping}

当系统参数$\boldsymbol{\theta}$未知时，将Backstepping与参数自适应律结合，在每个反推步骤中同时设计参数估计更新律。典型结构：状态方程为
\[
\dot{x}_i = f_i(\bar{x}_i) + \varphi_i^{\top}(\bar{x}_i)\boldsymbol{\theta} + g_i(\bar{x}_i)x_{i+1}
\]
Lyapunov函数中增加参数估计误差的二次项：
\[
V_i = V_{i-1} + \frac{1}{2}z_i^2 + \frac{1}{2}\tilde{\boldsymbol{\theta}}_i^{\top}\Gamma_i^{-1}\tilde{\boldsymbol{\theta}}_i
\]
其中$\tilde{\boldsymbol{\theta}}_i = \hat{\boldsymbol{\theta}}_i - \boldsymbol{\theta}$是参数估计误差，$\Gamma_i$是自适应增益矩阵。在每个步骤中适当选择虚拟控制律和参数更新律，保证$\dot{V}_i \le 0$。这种"调节函数"（tuning function）方法确保了最终的参数估计和状态跟踪同时收敛。

%%%%%%%%%%%%%%%%%%%%%%%%%%%%%%%%%%%%%%%%%%%%%%%%%%%%%%%%%%%%%%%%%%%%%%%%%%%%%%%%
\section{非线性观测器}

\subsection{高增益观测器}

对于可化为标准观测形式的非线性系统，可以设计高增益观测器。其基本思想是：利用足够大的观测器增益克服非线性不确定性。以SISO系统为例，若系统可写成
\begin{align*}
\dot{x}_1 &= x_2 + \varphi_1(x_1, u) \\
\dot{x}_2 &= x_3 + \varphi_2(x_1, x_2, u) \\
&\;\;\vdots \\
\dot{x}_n &= \varphi_n(\mathbf{x}, u) \\
y &= x_1
\end{align*}
则可构造高增益观测器：
\begin{align*}
\dot{\hat{x}}_1 &= \hat{x}_2 + \varphi_1(\hat{x}_1, u) + (\gamma_1/\varepsilon)(y - \hat{x}_1) \\
\dot{\hat{x}}_2 &= \hat{x}_3 + \varphi_2(\hat{x}_1, \hat{x}_2, u) + (\gamma_2/\varepsilon^2)(y - \hat{x}_1) \\
&\;\;\vdots \\
\dot{\hat{x}}_n &= \varphi_n(\hat{\mathbf{x}}, u) + (\gamma_n/\varepsilon^n)(y - \hat{x}_1)
\end{align*}
其中$\varepsilon > 0$充分小（高增益），系数$\gamma_i$使多项式$s^n + \gamma_1 s^{n-1} + \cdots + \gamma_n$为Hurwitz。$\varepsilon$越小，估计误差的收敛速度越快，但对测量噪声的敏感性也越强。

\subsection{滑模观测器}

将滑模技术应用于观测器设计，利用不连续的注入项使估计误差滑向零。滑模观测器的基本结构为
\[
\dot{\hat{\mathbf{x}}} = A\hat{\mathbf{x}} + \mathbf{\nu}
\]
其中$\mathbf{\nu}$是不连续的注入项，通常取$\mathbf{\nu} = \mathbf{L}\operatorname{sgn}(y - C\hat{\mathbf{x}})$。在滑模面上，等效注入等价于一个线性的Luenberger观测器，从而在滑模期间实现精确的状态重构。

滑模观测器的优势在于：
\begin{enumerate}
\item 对匹配不确定性和扰动具有完全鲁棒性（在滑动阶段）；
\item 可以实现有限时间状态估计（区别于渐近收敛）；
\item 在故障检测与隔离（FDI）中有广泛应用。
\end{enumerate}

%%%%%%%%%%%%%%%%%%%%%%%%%%%%%%%%%%%%%%%%%%%%%%%%%%%%%%%%%%%%%%%%%%%%%%%%%%%%%%%%
\section{Python实现：倒立摆的滑模控制仿真}

\begin{mycode}{滑模控制器实现}
\begin{lstlisting}[language=Python]
import numpy as np
import matplotlib.pyplot as plt

# 倒立摆参数
g, L, m = 9.81, 0.5, 0.2

def pend_dyn(x, u):
    """倒立摆非线性动力学 x = [theta, dtheta]"""
    theta, dtheta = x[0], x[1]
    ddtheta = (g/L) * np.sin(theta) + u / (m * L**2)
    return np.array([dtheta, ddtheta])

class SlidingModeController:
    def __init__(self, lam=5.0, eta=0.5, k=8.0, phi=0.05):
        self.lam = lam      # 滑模面参数 s = dtheta + lam*theta
        self.eta = eta      # 等速趋近律系数
        self.k = k          # 指数趋近律系数
        self.phi = phi      # 边界层厚度
    
    def sliding_surface(self, x):
        return x[1] + self.lam * x[0]
    
    def control(self, x):
        s = self.sliding_surface(x)
        theta, dtheta = x[0], x[1]
        
        # 等效控制（当 s = 0 且 \dot{s} = 0 时）
        u_eq = -m * L**2 * (self.lam * dtheta + (g/L) * np.sin(theta))
        
        # 趋近律（带边界层的饱和函数）
        b = m * L**2
        if abs(s) > self.phi:
            u_sw = -b * (self.k * np.sign(s) + self.eta * s)
        else:
            u_sw = -b * (self.k * s / self.phi + self.eta * s)
        
        return u_eq + u_sw

# 仿真
dt = 0.001
t = np.arange(0, 5, dt)
n = len(t)

smc = SlidingModeController(lam=5.0, eta=0.5, k=8.0, phi=0.05)

x = np.zeros((2, n))
u_arr = np.zeros(n)
s_arr = np.zeros(n)
x[:, 0] = [0.5, 0.0]  # 初始角30度

for i in range(n - 1):
    u_arr[i] = smc.control(x[:, i])
    s_arr[i] = smc.sliding_surface(x[:, i])
    xdot = pend_dyn(x[:, i], u_arr[i])
    x[:, i+1] = x[:, i] + xdot * dt

fig, axes = plt.subplots(3, 1, figsize=(10, 8))

axes[0].plot(t, x[0], 'b')
axes[0].set_ylabel(r'$\theta$ (rad)')
axes[0].grid(True)
axes[0].set_title('倒立摆滑模控制仿真')

axes[1].plot(t, x[1], 'r')
axes[1].set_ylabel(r'$\dot{\theta}$ (rad/s)')
axes[1].grid(True)

axes[2].plot(t, s_arr, 'g')
axes[2].axhline(y=0, color='k', linestyle='--', linewidth=0.8)
axes[2].set_ylabel('s (滑模面)')
axes[2].set_xlabel('Time [s]')
axes[2].grid(True)

plt.tight_layout()
plt.show()
\end{lstlisting}
\end{mycode}

\begin{mycode}{不同趋近律的对比}
\begin{lstlisting}[language=Python]
import numpy as np
import matplotlib.pyplot as plt

def simulate_smc(controller, x0, dt=0.001, T=5):
    n = int(T / dt)
    t = np.linspace(0, T, n)
    x = np.zeros((2, n))
    s = np.zeros(n)
    x[:, 0] = x0
    for i in range(n - 1):
        s[i] = controller.sliding_surface(x[:, i])
        u = controller.control(x[:, i])
        xdot = pend_dyn(x[:, i], u)
        x[:, i+1] = x[:, i] + xdot * dt
    return t, x, s

class SMC_Exponential:
    def __init__(self, lam=5.0, eta=0.5, k=8.0):
        self.lam, self.eta, self.k = lam, eta, k
    def sliding_surface(self, x):
        return x[1] + self.lam * x[0]
    def control(self, x):
        s = self.sliding_surface(x)
        u_eq = -m*L**2 * (self.lam*x[1] + (g/L)*np.sin(x[0]))
        u_sw = -m*L**2 * (self.k * np.sign(s) + self.eta * s)
        return u_eq + u_sw

class SMC_PowerRate:
    def __init__(self, lam=5.0, k=5.0, alpha=0.5):
        self.lam, self.k, self.alpha = lam, k, alpha
    def sliding_surface(self, x):
        return x[1] + self.lam * x[0]
    def control(self, x):
        s = self.sliding_surface(x)
        u_eq = -m*L**2 * (self.lam*x[1] + (g/L)*np.sin(x[0]))
        u_sw = -m*L**2 * self.k * np.abs(s)**self.alpha * np.sign(s)
        return u_eq + u_sw

# 对比仿真
smc_exp = SMC_Exponential()
smc_pwr = SMC_PowerRate()
t1, x1, s1 = simulate_smc(smc_exp, [0.5, 0.0])
t2, x2, s2 = simulate_smc(smc_pwr, [0.5, 0.0])

fig, axes = plt.subplots(2, 2, figsize=(12, 8))
axes[0, 0].plot(t1, x1[0], label='Exponential')
axes[0, 0].plot(t2, x2[0], '--', label='Power Rate')
axes[0, 0].set_ylabel(r'$\theta$ (rad)')
axes[0, 0].legend()
axes[0, 0].grid(True)

axes[0, 1].plot(t1, s1, label='Exponential')
axes[0, 1].plot(t2, s2, '--', label='Power Rate')
axes[0, 1].axhline(y=0, color='k', linestyle=':', lw=0.8)
axes[0, 1].set_ylabel('s (滑模面)')
axes[0, 1].legend()
axes[0, 1].grid(True)

axes[1, 0].plot(x1[0], x1[1], 'b')
axes[1, 0].plot([-0.5, 0.5], [2.5, -2.5], 'r--', lw=1, label='s=0')
axes[1, 0].set_xlabel(r'$\theta$')
axes[1, 0].set_ylabel(r'$\dot{\theta}$')
axes[1, 0].set_title('相平面（指数趋近律）')
axes[1, 0].legend()
axes[1, 0].grid(True)

axes[1, 1].plot(x2[0], x2[1], 'g')
axes[1, 1].plot([-0.5, 0.5], [2.5, -2.5], 'r--', lw=1, label='s=0')
axes[1, 1].set_xlabel(r'$\theta$')
axes[1, 1].set_ylabel(r'$\dot{\theta}$')
axes[1, 1].set_title('相平面（幂次趋近律）')
axes[1, 1].legend()
axes[1, 1].grid(True)

plt.tight_layout()
plt.show()
\end{lstlisting}
\end{mycode}

%%%%%%%%%%%%%%%%%%%%%%%%%%%%%%%%%%%%%%%%%%%%%%%%%%%%%%%%%%%%%%%%%%%%%%%%%%%%%%%%
\section{小结}

非线性控制是一个广阔而活跃的研究领域。本章介绍了其基本框架和核心方法：

\begin{enumerate}
\item \textbf{非线性特性}：多平衡点、极限环、混沌和分岔使非线性系统远较线性系统复杂；
\item \textbf{Lyapunov稳定性理论}：以"能量函数"的概念为核心，是分析非线性系统稳定性的最基本工具；Lasalle不变原理推广了Lyapunov方法的应用范围；
\item \textbf{滑模控制}：利用不连续控制将状态驱动到预设滑模面，对匹配扰动具有完全鲁棒性，但需注意抖振问题；
\item \textbf{反馈线性化}：利用微分几何工具通过坐标变换和反馈将非线性系统精确转化为线性系统；Lie导数、Lie括号和相对阶是其中的核心概念；
\item \textbf{Backstepping}：针对严格反馈系统的递归Lyapunov设计方法，可灵活处理非线性和参数不确定性；
\item \textbf{非线性观测器}：高增益观测器和滑模观测器为非线性系统的状态估计提供了有效方案。
\end{enumerate}

现代非线性控制理论仍在持续发展，包括事件触发控制、预设性能控制、基于障碍Lyapunov函数的约束控制等新方向。下一章将进一步介绍模型预测控制、鲁棒控制和智能控制等高级方法。
